\chapter{Dalla sintassi astratta alla grammatica}

Nel capitolo precedente è stata introdotta la sintassi astratta del linguaggio Racing Choreographies insieme alla semantica operazionale.

Tuttavia, la sintassi astratta descrive la struttura logica dei programmi indipendentemente dalla loro rappresentazione testuale e per poter scrivere ed eseguire programmi è necessario definire una sintassi concreta e una grammatica per il parser.

Il processo seguito è stato quello di progettare una sintassi concreta poi formalizzata in una grammatica espressa in Extended Backus--Naur Form (EBNF) e tradotta infine in una grammatica ANTLR \cite{antlr4tool}.

Ai fini della comprensione del linguaggio viene illustrato il passo più rilevante ovvero la progettazione della sintassi concreta, che stabilisce come i costrutti teorici presentati nel capitolo precedente vengono effettivamente scritti nei programmi.

La grammatica completa utilizzata per l’implementazione del parser è disponibile nella repository del progetto nel file \texttt{RacingChoreo.g4} nella cartella \texttt{/grammar} \cite{rcparserrepo}.

\section{Principi e requisiti della sintassi concreta}

La progettazione della sintassi concreta di Racing Choreographies è stata guidata dalle seguenti scelte progettuali secondo uno stile C-like.

\subsection*{Separazione tra definizioni e coreografia principale}

Un programma è composto da zero o più definizioni di procedura seguite da una coreografia principale identificata dalla keyword \texttt{main}. Questa scelta rende esplicito il punto di ingresso dell’esecuzione del programma.

\subsection*{Uso di keyword esplicite}

I costrutti principali del linguaggio sono introdotti tramite keyword riservate, tra cui \texttt{proc}, \texttt{main}, \texttt{call}, \texttt{race}, \texttt{discharge}, \texttt{if} ed \texttt{else}.

L’uso di parole chiave esplicite facilita la distinzione tra identificatori del programma e costrutti del linguaggio.

\subsection*{Struttura a blocchi}

Le coreografie sono racchiuse tra parentesi graffe \texttt{\{ ... \}}, come accade in molti linguaggi di programmazione imperativi in quanto questo consente di organizzare in modo chiaro sequenze di istruzioni e rami condizionali.

\subsection*{Terminazione implicita della coreografia}

Il termine di terminazione \texttt{0}, presente nella sintassi astratta, non viene rappresentato esplicitamente nella sintassi concreta. La chiusura di un blocco corrisponde implicitamente alla terminazione della coreografia contenuta.

\subsection*{Sequenzialità esplicita tramite punto e virgola}

Le istruzioni e le chiamate di procedura sono terminate dal punto e virgola \texttt{;}, rendendo esplicita la sequenza delle azioni. I costrutti condizionali, invece, sono strutturati come blocchi e non richiedono il punto e virgola finale.


\subsection*{Distinzione tra condizionali locali e su race}

Le due forme di condizionale previste dalla sintassi astratta sono rese sintatticamente distinte:

\begin{center}
\texttt{if (p.e) \{ ... \} else \{ ... \}}
\end{center}

\begin{center}
\texttt{if (s[k]) \{ ... \} else \{ ... \}}
\end{center}

Nel primo caso la condizione dipende dalla valutazione locale di un’espressione appartenente a un processo. Nel secondo caso la condizione dipende dall’esito di una race identificata da \texttt{s[k]}.

\subsection*{Espressioni semplici}

Le espressioni del linguaggio sono volutamente limitate a valori letterali o identificatori. Non sono presenti operatori aritmetici poiché il focus del linguaggio è descrivere interazioni tra processi piuttosto che calcoli locali.
\section{Sintassi concreta attraverso un esempio}
Per illustrare la sintassi concreta del linguaggio Racing Choreographies si riprende l'esempio completo di programma introdotto precedentemente che utilizza i costrutti del linguaggio.
Il programma descrive una semplice interazione tra quattro processi: un client \texttt{c}, due worker \texttt{w1} e \texttt{w2}, e un server \texttt{s}. Il client invia una richiesta al server; i due worker ricevono la richiesta e competono per fornire una risposta tramite una \texttt{race}. Il valore prodotto dal worker vincitore viene inviato al client, mentre il valore del worker perdente viene successivamente recuperato tramite il costrutto di \texttt{discharge}.
Così è come appare il programma completo secondo la sintassi concreta:
\begin{lstlisting}
proc Request(c, w1, w2, s) {

    c.req -> s.req;

    w1.r = req;
    w2.r = req;

    race s[k] : w1.r , w2.r -> s.ans;

    if (s[k]) {
        s.ans -> c.res;
        s -> c[FromW1];
        discharge s[k] : w2 -> s.lost;
    } else {
        s.ans -> c.res;
        s -> c[FromW2];
        discharge s[k] : w1 -> s.lost;
    }
}

main {
    call Request(c, w1, w2, s);
}
\end{lstlisting}

Il programma è composto da una definizione di procedura e da una coreografia principale.

La procedura \texttt{Request} descrive l’interazione tra i quattro processi coinvolti, mentre il blocco \texttt{main} rappresenta il punto di ingresso dell’esecuzione del programma e invoca la procedura tramite la keyword \texttt{call}.

La prima istruzione

\begin{tcolorbox}[colback=white,colframe=black!35,boxrule=0.4pt]
\begin{lstlisting}
c.req -> s.req;
\end{lstlisting}
\end{tcolorbox}

rappresenta una comunicazione globale: il processo \texttt{c} invia al processo \texttt{s} il valore della variabile \texttt{req}, che viene memorizzato nella variabile \texttt{req} del server.

Le istruzioni successive

\begin{tcolorbox}[colback=white,colframe=black!35,boxrule=0.4pt]
\begin{lstlisting}
w1.r = req;
w2.r = req;
\end{lstlisting}
\end{tcolorbox}

sono assegnamenti locali, che permettono ai worker di preparare il valore che verrà utilizzato nella competizione.

La competizione è espressa dal costrutto di \texttt{race}:

\begin{tcolorbox}[colback=white,colframe=black!35,boxrule=0.4pt]
\begin{lstlisting}
race s[k] : w1.r , w2.r -> s.ans;
\end{lstlisting}
\end{tcolorbox}

In questa istruzione i processi \texttt{w1} e \texttt{w2} competono per fornire un valore al processo \texttt{s}. Il valore del processo vincitore viene assegnato alla variabile \texttt{s.ans}. L’identificatore \texttt{s[k]} rappresenta la race con id \texttt{k} e consente di riferirsi successivamente al suo esito.

L’esito della competizione viene utilizzato nel costrutto condizionale

\begin{tcolorbox}[colback=white,colframe=black!35,boxrule=0.4pt]
\begin{lstlisting}
if (s[k]) { ... } else { ... }
\end{lstlisting}
\end{tcolorbox}

che seleziona il ramo di esecuzione in base al vincitore della race.

Nel ramo selezionato il valore vincente viene inviato al client tramite una comunicazione:

\begin{tcolorbox}[colback=white,colframe=black!35,boxrule=0.4pt]
\begin{lstlisting}
s.ans -> c.res;
\end{lstlisting}
\end{tcolorbox}

Successivamente il server invia al client una selezione di etichetta (\texttt{FromW1} oppure \texttt{FromW2}) per indicare quale worker ha prodotto il risultato.

Infine, il costrutto

\begin{tcolorbox}[colback=white,colframe=black!35,boxrule=0.4pt]
\begin{lstlisting}
discharge s[k] : w1 -> s.lost
discharge s[k] : w2 -> s.lost
\end{lstlisting}
\end{tcolorbox}

consente di recuperare il valore prodotto dal processo perdente della race. Questo valore viene memorizzato nella variabile \texttt{s.lost}, rendendo esplicita la sincronizzazione con il processo che non ha vinto la competizione.